\sectioncentered*{Реферат}
\thispagestyle{empty}

\MakeUppercase{наименование темы дипломного проекта}: персональный ассистент организации времени обучающихся/ Е.~В. Веркович. -- Минск : БГУИР, \the\year{2026}, -- п.з. -- ~\pageref*{LastPage}~с., чертежей (плакатов) -- 6 л. формата А1.

\vspace{4\parsep}

%Дипломный проект выполнен на 6 листах А1 с пояснительной запиской на на~\pageref*{LastPage} страницах, без приложений справочного или информационного характера. Пояснительная записка включает 4 раздела, 27 рисунков, 3 таблицы, 9 формул, 29 литературных источников.

Целью дипломного проекта является оптимизация организации времени обучающихся за счет делегирования планирования учебных задач интеллектуальному ассистенту.

Предметом исследования являются методы интеллектуальной поддержки организации времени обучающихся, а также архитектурные и программные подходы к реализации персональных ассистентов, ориентированных на декомпозицию учебных задач и планирование индивидуальной нагрузки.

Разработанный программный продукт позволяет формировать персонализированный черновик расписания на основе текстового запроса пользователя, выполнять декомпозицию долгосрочных учебных задач на подзадачи, распределять их по свободным временным интервалам и предоставлять пользователю средства подтверждения, отклонения и ручной корректировки предложенного плана.

В первом разделе пояснительной записки проведен анализ предметной области, рассмотрены психолого-педагогические причины академической прокрастинации, исследованы существующие аналоги и современные подходы к разработке интеллектуальных ассистентов организации времени.

Во втором разделе выполнено проектирование системы: определены функциональные и нефункциональные требования, разработаны варианты использования, логическая архитектура, модель данных, а также диаграммы контроллеров, сервисов и репозиториев.

В третьем разделе описана реализация программного комплекса, обоснован выбор технических средств, рассмотрены алгоритм размещения подзадач по временным слотам, обработка структурированного вывода языковой модели и жизненный цикл черновика расписания, а также приведена демонстрация работоспособности системы.

В четвертом разделе приведено технико-экономическое обоснование разрабатываемого программного продукта и его использования по индивидуальному заказу.

Результатом дипломного проектирования является программный комплекс персонального ассистента организации времени обучающихся, включающий серверную часть, агентный модуль, подсистему хранения данных и клиентский интерфейс с календарной визуализацией черновика расписания.

\clearpage